\notepage{60}{Position-spin wavefunctions and momentum}
For a product state $\ket{\Psi}=\ket{\varphi}\otimes\ket{\chi}$,
the $S_z$ representation first replaces the spin factor by its two
coefficients:
\[
 \begin{bmatrix}
  \braket{+}{\chi}\ket{\varphi}\\
  \braket{-}{\chi}\ket{\varphi}
 \end{bmatrix}.
\]
Taking spatial position matrix elements gives the position-plus-$S_z$
representation,
\begin{align*}
 \boldsymbol\Psi(\mathbf r)
 &=\begin{bmatrix}
   \braket{+}{\chi}\braket{\mathbf r}{\varphi}\\
   \braket{-}{\chi}\braket{\mathbf r}{\varphi}
  \end{bmatrix}
 =\begin{bmatrix}\Psi_+(\mathbf r)\\\Psi_-(\mathbf r)\end{bmatrix}\\
 &=\Psi_+(\mathbf r)\begin{bmatrix}1\\0\end{bmatrix}
  +\Psi_-(\mathbf r)\begin{bmatrix}0\\1\end{bmatrix}.
\end{align*}
The two-component expression also represents general, nonfactorized
states when the two functions are allowed to vary independently.

\smallheading{Momentum acts on the spatial factor}
In the product Hilbert space,
\[
 \mathbf P_{r+s}=\mathbf P_r\otimes I_s.
\]
Therefore its action in the position-plus-$S_z$ representation is
\[
 \mathbf P_{r+s}\ \longleftrightarrow\
 \begin{bmatrix}
  -\ii\hbar\boldsymbol\nabla&0\\
  0&-\ii\hbar\boldsymbol\nabla
 \end{bmatrix},
\]
so
\[
 (\mathbf P\boldsymbol\Psi)(\mathbf r)=
 \begin{bmatrix}
  -\ii\hbar\boldsymbol\nabla\Psi_+(\mathbf r)\\
  -\ii\hbar\boldsymbol\nabla\Psi_-(\mathbf r)
 \end{bmatrix}.
\]
The gradient does not mix spin-up and spin-down components.
